\section{Paris--Harrington principle}

\begin{definition}[Relatively large homogeneous set]
\label{def:ph-large}
\lean{ParisHarringtonPrinciple.RelativelyLarge,ParisHarringtonPrinciple.IsHomogeneous}
\leanok
The blueprint uses the standard notions of a relatively large subset of $[1,N]$ and a homogeneous set for a coloring of $k$-subsets.
\end{definition}

\begin{theorem}[Infinite Ramsey theorem]
\label{thm:ph-infinite-ramsey}
\lean{ParisHarringtonPrinciple.infiniteRamsey}
\leanok
Every finite coloring of the $k$-subsets of an infinite set has an infinite homogeneous subset.
\end{theorem}

\begin{lemma}[Ultrafilter fiber]
\label{lem:ph-ultrafilter}
\lean{ParisHarringtonPrinciple.ultrafilter_fiber}
\leanok
For a map into a finite color set, some color fiber belongs to any ultrafilter.
\end{lemma}

\begin{lemma}[Finite subset of prescribed cardinality]
\label{lem:ph-finset}
\lean{ParisHarringtonPrinciple.exists_finset_card_eq}
\leanok
An infinite subset of $\mathbb N$ contains a finite subset of every prescribed finite cardinality.
\end{lemma}

\begin{theorem}[Paris--Harrington]
\label{thm:paris-harrington}
\lean{ParisHarringtonPrinciple.MainTheorem}
\leanok
\uses{def:ph-large,thm:ph-infinite-ramsey,lem:ph-ultrafilter,lem:ph-finset}
For $0<r$ and $k\le m$, there is an $N$ such that every $r$-coloring of the $k$-subsets of $[1,N]$ has a relatively large homogeneous set.
\end{theorem}
